\section{问题重述与分析}
\label{sec:intro}

\subsection{问题背景}
大语言模型的训练可类比为培养一名学生：参数量 $N$ 相当于脑容量，训练数据量 $D$ 相当于阅读量，计算量 $C$ 相当于教育经费，数据质量 $Q$ 相当于教材质量，领域配比 $p$ 相当于各科课时分配。经典标度律刻画损失随规模变化的经验关系，但数据墙（Data Wall）背景下，高质量语料日益稀缺，数据质量与领域配比对模型能力的作用已不可忽视——小参数模型凭借高质量数据在引文准确率等任务上超过大模型的公开案例，正是本题的现实动因。

本题要求在算力有限的条件下，建立参数规模、数据规模、数据质量、领域配比与交叉熵损失之间的定量关系，并研究算力约束下的资源配置问题。四问递进：问题一刻画数据质量与配比；问题二将质量与配比纳入广义标度律；问题三在算力预算下求解多维资源联合优化并识别结构性转移；问题四分离规模扩张与非规模技术进步的贡献，预测开源模型能力前沿。

\subsection{数据概况}
全部数据来自题目附件 real\_attachments/，以真实公开数据为主、真实数据校准的半合成数据为补充。附件 A（问题一）含 SlimPajama-Meta-rater 质量信号（A1 抽样集 51,230 条，A2/A3 为 arxiv、github 扩展集）与 RegMix 配方实验（A4--A15，512 组训练配方及三尺度检验/外推表）；附件 B（问题二）含 Pythia 真实训练轨迹（B1）、Cerebras 半合成轨迹（B2）、插值轨迹（B3）、跨族与文献基准（B4/B5）及唯一含质量分的 NQ 半合成实验（B6--B8）；附件 C（问题三、四）含 Open LLM Leaderboard 六维评测（C1/C2）、时序（C3）、宏观元数据（C4）、Loss--Benchmark 桥接（C5/C6，按可比性分层）、模型架构与上下文长度元数据（C7）及逐任务评测 JSON（C8，必用）。各文件编号、字段与性质的完整说明以题目《数据说明》为准；凡属半合成、外推或估算性质的数据，本文在使用处显式标注其证据等级。

% 问题一问题重述
% 本文件只描述问题背景与题目要求，不包含问题一的模型、结果或结论。

\subsection{问题一问题重述}
\label{sec:q1-restatement}

\subsubsection{问题背景}

\paragraph{大语言模型训练中的数据质量问题}
随着大语言模型规模不断扩大，模型性能不仅受到参数量和训练数据量的影响，也与训练数据本身的质量密切相关。在相同的模型规模和训练预算下，高质量、低重复、覆盖合理且具有良好代表性的语料，往往能够带来更低的训练损失和更好的模型表现。

然而，训练语料通常来源复杂，包含网页文本、论文、代码、书籍及其他类型的数据。不同语料在重复程度、语言质量、信息密度、毒性、可读性和领域覆盖等方面存在明显差异。因此，如何对多来源语料进行统一、客观且可比较的质量评价，是大语言模型数据治理和训练优化中的重要问题。

\paragraph{多指标质量评价与质量冲突}
题目给出了多类反映语料质量的指标，其中既包括可以直接计算的标量指标，也包括需要进一步汇总处理的列表型指标。不同指标的量纲、取值范围和优劣方向并不一致，有些指标数值越大表示质量越高，而另一些指标数值越小反而表示质量越好。

此外，不同质量指标之间可能存在相互矛盾的情况。例如，某一批数据可能具有较高的信息密度，但同时存在较多重复内容；也可能在语言可读性方面表现较好，却在领域覆盖或安全性方面存在不足。若简单地将各项指标加权为一个总分，可能掩盖数据的结构性缺陷。因此，需要建立综合评价模型，并进一步识别和处理指标之间的质量冲突。

\paragraph{领域配比与模型训练损失}
训练数据不仅具有整体质量差异，还具有领域组成上的差异。不同领域的数据在训练语料中的占比会影响模型对各类知识和任务的学习效果，进而影响模型的交叉熵损失。

题目提供了多组不同领域配比及其对应的训练损失数据。设训练语料包含 $17$ 个领域，领域配比记为
\begin{equation}
  p=(p_1,p_2,\ldots,p_{17}),
  \qquad p_i\geq 0,\quad \sum_{i=1}^{17}p_i=1.
  \label{eq:q1-restatement-simplex}
\end{equation}
因此，领域配比建模是在单纯形约束下研究各领域占比与交叉熵损失之间的关系，并据此分析不同领域组合对模型训练效果的影响。

\subsubsection{问题提出}

问题一围绕训练数据的质量评价、质量冲突和领域配比三个方面展开，具体包括以下三项任务。

\paragraph{第一小问：数据质量评价}
利用题目提供的多维质量指标，统一各指标的方向、量纲和取值范围，建立综合质量评价模型，得到样本级、语料级和领域级的数据质量评分。模型需要充分利用题目给出的质量信息，对抽样数据和扩展数据进行评价，分析不同数据来源及不同语料领域之间的质量差异，并形成可用于后续建模的质量评分 $Q$。

\paragraph{第二小问：质量冲突消解}
在综合质量评价的基础上，定义数据质量冲突的判定标准，识别不同质量指标之间的不一致关系，并分析产生冲突的可能原因。针对存在冲突的数据，建立相应的冲突消解方法和修正后的综合评价模型，考察冲突处理前后数据质量排序及评价结果的变化，并在扩展数据上检验该方法的稳定性和适用性。

\paragraph{第三小问：领域配比建模}
利用题目给出的领域配比和交叉熵损失数据，建立 $17$ 个领域配比 $p$ 与模型训练损失 $\mathrm{Loss}$ 之间的定量关系。模型需要在满足式~\eqref{eq:q1-restatement-simplex} 的条件下，分析不同领域的影响及其组合效应，并利用训练数据建立关系模型，再通过检验数据考察模型的预测能力和稳健性。对于更大规模模型对应的数据，还需要讨论模型在外推场景下的适用范围，并说明部分领域缺少直接损失观测时的处理方式。

通过上述三项任务，问题一形成从“多指标质量测度”到“质量冲突处理”，再到“领域配比与训练损失关系分析”的建模链条，为后续研究数据质量、模型规模和算力资源之间的关系提供质量评分 $Q$ 与领域配比 $p$ 等基础变量。


\subsection{建模任务的拆解与本文工作}
\begin{itemize}
  \item \textbf{问题一}（第~\ref{sec:q1}~节）：22 维质量指标经方向统一与稳健归一化后，以 CRITIC 客观赋权与加权距离 TOPSIS 输出样本级质量分，聚合至领域级与语料级；对同一文本上的指标不一致定义冲突并构造有限补偿消解模型；对 17 域配比建立单纯形约束下的同尺度响应面模型。
  \item \textbf{问题二}：以 B1 为主拟合数据建立含质量与配比的广义标度律，B2--B5 分层验证，B6--B8 引入质量维度；分析边际效用与弹性，推导质量—规模可替代条件。（本文版本给出接口与探索性结果，正式联合拟合以问题一输出的 $Q_g$ 为输入。）
  \item \textbf{问题三}：在 $C=6ND+\Delta C_Q+\Delta C_S$ 的算力结构下，对低/中/高三档预算求解 $N,D,Q,p$ 联合配置，给出结构性转移的数学定义与识别方法，并依据 C7 对上下文长度 $S$ 做外生敏感性分析。
  \item \textbf{问题四}：建立动力学/因果分解模型分离规模扩张与非规模技术进步贡献；经 C6 桥接把损失映射为 Benchmark 得分；在算力增长放缓情景下预测未来 12--24 个月开源能力前沿并给出不确定性分析。
\end{itemize}
